\section{Stable unfolded probabilities}\label{sec:stability}

We now compute the survival probabilities in the independent unfolding.
Throughout this section, $G$ is any graph satisfying
Proposition~\ref{prop:prefix}; every estimate is uniform over such graphs.
Spencer's branching model has survival probability $(1+Qc)^{-1/Q}$
\cite[Section~3, equation~(9)]{Spencer1995}. With $Q=2$ and $c=Dt$, its
scalar comparison function is
\begin{equation}\label{eq:scalar-survival}
 q(t)=(1+2Dt)^{-1/2},\qquad 0\le t\le1.
\end{equation}
It solves $q'(t)=-Dq(t)^3$ with $q(0)=1$.  We will show that the
probabilities $q_e(t)$ and $q_{e,T}(t)$ from Section~\ref{sec:queries}
are uniformly asymptotic to $q(t)$.  The essential issue is stability:
small discrepancies between edge types must remain small throughout
the time interval.

\subsection{The exact equation and its linearization}

Use the increasing change of variables
\[
 s=\tfrac12\log(1+2Dt),\qquad
 S=\tfrac12\log(1+2D),\qquad
 t(s)=\frac{e^{2s}-1}{2D}.
\]
Thus $0\le s\le S$, $q(t(s))=e^{-s}$, and $S\le\log n$ for
all sufficiently large $n$.  Since $0\le I_{e,T}(t)\le t$, the product
identities in Proposition~\ref{prop:products} give $q_e(t)>0$ for
$t<1$.  On any interval on which all product factors remain positive,
define the vector $z(s)\in\mathbb R^{\cE}$ by
\[
 z_e(s)=\log\frac{q_e(t(s))}{q(t(s))}.
\]
In particular, $z(0)=0$.  Positivity at $t=1$ will follow from the
estimates below.

Let $A$ be the real symmetric matrix indexed by $\cE$, with
$A_{ef}=1$ if $e\ne f$ belong to a common triangle of $G$, and
$A_{ef}=0$ otherwise.  Each row indexed by $e$ has sum $2d_e$.
Linearity of the triangle hypergraph ensures that each pair of distinct
edges contributes at most once to $A$.

Differentiating the products is legitimate by
Proposition~\ref{prop:products}.  For $T=\{e,f,g\}$ we have
\[
 I'_{e,T}(t)=q_{f,T}(t)q_{g,T}(t),\qquad
 q_{f,T}(t)=\frac{q_f(t)}{1-I_{f,T}(t)},\qquad
 \frac{dt}{ds}=\frac1{Dq(t)^2}.
\]
Here and below a prime on $I$ denotes differentiation in $t$, and a
prime on $z$ denotes differentiation in $s$.  Taking the logarithmic
derivative of $q_e$ consequently gives the exact identity
\begin{equation}\label{eq:exact-z}
 z'_e(s)=1-\frac1D
 \sum_{\substack{T\in\cT\\T=\{e,f,g\}}}
 \frac{\exp(z_f(s)+z_g(s))}
 {(1-I_{e,T}(t))(1-I_{f,T}(t))(1-I_{g,T}(t))},
 \qquad t=t(s).
\end{equation}
The factor indexed by $e$ comes from differentiating its product;
the other two factors come from the two omitted-parent probabilities.

For vectors, $\lVert\cdot\rVert_\infty$ denotes the maximum norm and
$\lVert\cdot\rVert_2$ the Euclidean norm.  For matrices we use the
corresponding induced operator norms.  Fix
\begin{equation}\label{eq:bootstrap-exponent}
 0<a<\min\{c,2\epsilon\},\qquad \delta_n=n^{-a}.
\end{equation}
Suppose, provisionally, that on an interval starting at zero,
\begin{equation}\label{eq:bootstrap-assumptions}
 \lVert z(s)\rVert_\infty\le\delta_n,
 \qquad I_{e,T}(t(s))\le\tfrac12
 \quad(e\in T\in\cT).
\end{equation}
The omitted-parent identity then implies
$q_{e,T}(t)\le2e^{\delta_n}q(t)$.  Integrating yields, on the same
interval,
\begin{equation}\label{eq:I-bootstrap}
 I_{e,T}(t)
 \le4e^{2\delta_n}\int_0^t\frac{dv}{1+2Dv}
 =\frac{2e^{2\delta_n}}D\log(1+2Dt)
 =O\!\left(\frac{\log n}{D}\right).
\end{equation}
In particular, all three denominator factors in~\eqref{eq:exact-z}
are $1+O(\log n/D)$.  Expanding the exponential and reciprocal
factors gives, uniformly for each summand,
\[
 \frac{e^{z_f+z_g}}
 {(1-I_{e,T})(1-I_{f,T})(1-I_{g,T})}
 =1+z_f+z_g+
 O\!\left(\delta_n^2+\frac{\log n}{D}\right).
\]
The sum of the linear terms is exactly $(Az)_e$.  Since
$d_e/D=1+O(n^{-c})$, we obtain the conditional estimate
\begin{equation}\label{eq:linearized-error}
 \left\lVert z'(s)+(A/D)z(s)\right\rVert_\infty
 \le C\left(n^{-2a}+n^{-c}+\frac{\log n}{D}\right).
\end{equation}
The constant $C$ is independent of the graph and the provisional
interval.  The general bound
$\lVert e^{-sA/D}\rVert_\infty\le
e^{s\lVert A/D\rVert_\infty}$ gives a loss of order $D$ at $s=S$,
too large to guarantee control of the prefix error $n^{-c}$.
We will instead prove a polylogarithmic bound, sufficient to turn
\eqref{eq:linearized-error} into a strict improvement
of~\eqref{eq:bootstrap-assumptions}.

\subsection{From link spectra to a uniform semigroup bound}

The passage from even closed-walk counts to spectral control follows
the sparse quasirandom trace method; compare Chung and
Graham~\cite[Theorem~1 and Facts~7--8]{ChungGraham2002}.
We prove the quantitative form required by the link hypotheses and then
transfer it to the maximum norm on the full edge space.

\begin{proposition}\label{prop:semigroup}
There are fixed constants $C_1,J<\infty$ such that every graph $G$
satisfying the degree, codegree and link-cycle conditions of
Proposition~\ref{prop:prefix} obeys
\begin{equation}\label{eq:semigroup}
 \sup_{0\le s\le\log n}
 \left\lVert e^{-sA/D}\right\rVert_\infty
 \le C_1(1+\log n)^J
\end{equation}
for all sufficiently large $n$.
\end{proposition}

\begin{proof}
We first approximate the adjacency matrix on each star by an averaging
operator.  We then combine these averages and estimate the effect of
the error.  All constants in the proof are uniform in $G$.

\paragraph{Closed walks in a link.}
For $u\in[n]$, let $\cE_u=\{uv:v\in N_G(u)\}$ and
$N_u=|\cE_u|$.  Identify $uv\in\cE_u$ with $v\in N_G(u)$, and let
$M_u$ be the adjacency matrix of the link $G[N_G(u)]$ on this index
set.  We have $N_u=(1\pm n^{-c})np$, and every row sum of $M_u$ is
$(1\pm n^{-c})D$, by the codegree assumption.  Put $r=r_0=8000$.

The trace $\operatorname{tr}(M_u^r)$ counts closed walks of length
$r$ with a specified starting position.  Walks whose $r$ positions
are distinct are precisely embeddings of the labeled cycle $C_r$;
their number is $(1\pm n^{-c})D^r$.  For all other walks, classify
the patterns of identifications among the $r$ positions.  There are
only finitely many patterns, since $r$ is fixed.  The support of any
realized pattern is a connected simple graph; patterns forcing a loop
have no realizations.

If the support is a tree, every support edge is traversed at least
twice: a walk crossing the cut obtained by deleting that edge must
cross back.  Thus the support has at most $r/2$ edges.  Choose the
image of one vertex and then expose the other vertices along this
tree.  The link has maximum degree at most $2D$, so this gives at most
$O(np)(2D)^{r/2}$ realizations of each pattern.

If the support contains a cycle, let it have $j<r$ vertices and choose
one of its simple cycles, of length $b\ge3$.  The link-cycle condition
gives $O(D^b)$ choices for this labeled cycle.  A forest rooted at its
vertices attaches all the other vertices of the connected support,
using at most $(2D)^{j-b}$ choices.  Ignoring its remaining edges only
increases the count.  The resulting bound is $O(D^j)$, and hence
$O(D^{r-1})$ for each such pattern.  These cases exhaust the
non-injective walks.  Dividing their total contribution by $D^r$
shows that
\begin{equation}\label{eq:link-trace}
 \begin{aligned}
 \operatorname{tr}\bigl((M_u/D)^r\bigr)
 &=1+O\!\left(n^{-c}+D^{-1}+npD^{-r/2}\right)\\
 &=1+O(n^{-c_*}),\qquad c_*:=\min\{c,1/1000\}.
 \end{aligned}
\end{equation}
Indeed, $D\sim n^{1/1000}$ and
$npD^{-4000}\sim n^{-6999/2000}$.

\paragraph{The leading eigendirection.}
Write $B_u=M_u/D$, and order its real eigenvalues as
$\lambda_1\ge\lambda_2\ge\cdots$.  The row-sum bound gives
$\lambda_1\le1+n^{-c}$; the Rayleigh quotient of the normalized
all-ones vector gives $\lambda_1\ge1-n^{-c}$.  Because $r$ is even,
\[
 0\le\sum_{j\ge2}\lambda_j^r
 =\operatorname{tr}(B_u^r)-\lambda_1^r
 =O(n^{-c_*}).
\]
Consequently, with the fixed positive constant
$\gamma=c_*/r$, all eigenvalues other than $\lambda_1$ have absolute
value $O(n^{-\gamma})$.  In particular, $\lambda_1$ is simple for
large $n$.

Let $\mathbf j_u=N_u^{-1/2}\mathbf1$ and let $v_u$ be a unit
eigenvector for $\lambda_1$.  The row sums also give
$\lVert(B_u-\mathrm{Id})\mathbf j_u\rVert_2\le n^{-c}$.
Write $\mathbf j_u=\alpha v_u+w$, where $w\perp v_u$.  Every
eigenvalue of $B_u-\mathrm{Id}$ on $v_u^\perp$ has absolute value at
least $1/2$ for large $n$, so $\lVert w\rVert_2\le2n^{-c}$.
The orthogonal projectors onto the lines spanned by $v_u$ and
$\mathbf j_u$ therefore differ in operator norm by at most
$2n^{-c}$.  Combining this with the eigenvalue estimates gives
\begin{equation}\label{eq:star-approximation}
 \left\lVert \frac{M_u}{D}-\frac{\mathbf1\mathbf1^{\mathsf T}}{N_u}
 \right\rVert_2=O(n^{-\gamma}).
\end{equation}

\paragraph{Summing the star averages.}
Extend every star matrix by zero to the full edge space
$\mathbb R^{\cE}$, and let $P$ be the sum of the averaging matrices
$\mathbf1\mathbf1^{\mathsf T}/N_u$ on the stars.  Distinct edges
lying in a common triangle have a unique common endpoint, so $A$ is
the sum of the extended matrices $M_u$.  Define
\[
 R=A/D-P.
\]
The matrix $P$ is symmetric and positive semidefinite.  For every
$x\in\mathbb R^{\cE}$, estimate~\eqref{eq:star-approximation} gives
\[
 |x^{\mathsf T}Rx|
 \le Cn^{-\gamma}\sum_{u\in[n]}
          \lVert x|_{\cE_u}\rVert_2^2
 =2Cn^{-\gamma}\lVert x\rVert_2^2.
\]
Each edge lies in exactly two stars, which explains the last equality
and prevents a loss proportional to the number of vertices.  Since
$R$ is symmetric, it follows that $\lVert R\rVert_2=O(n^{-\gamma})$.
Every row of $P$ has sum two, while a row of $A/D$ has sum $2d_e/D$.
Both matrices are entrywise nonnegative, so their absolute row sums
also give
\begin{equation}\label{eq:R-norms}
 \lVert R\rVert_2=O(n^{-\gamma}),\qquad
 \lVert R\rVert_\infty=O(1).
\end{equation}

We have thus obtained a small Euclidean error, whereas the evolution
in~\eqref{eq:linearized-error} requires the maximum norm.  The special
form of $P$ supplies this stronger control.  Define
\[
 \begin{aligned}
 U:\mathbb R^{[n]}&\longrightarrow\mathbb R^{\cE},
 & (Uy)_{uv}&=y_u+y_v,\\
 V:\mathbb R^{\cE}&\longrightarrow\mathbb R^{[n]},
 & (Vx)_u&=\frac1{N_u}\sum_{e\in\cE_u}x_e.
 \end{aligned}
\]
Then $P=UV$, $\lVert U\rVert_\infty=2$, and
$\lVert V\rVert_\infty=1$.  If $W$ is the neighbor random-walk
matrix, defined by
$(Wy)_u=N_u^{-1}\sum_{v\in N_G(u)}y_v$, then
$VU=\mathrm{Id}+W$ and $\lVert W\rVert_\infty=1$.
The identity $(UV)^j=U(VU)^{j-1}V$ for $j\ge1$, applied to the
exponential power series, yields
\[
 E_0(s):=e^{-sP}
 =\mathrm{Id}-U\int_0^s e^{-v(\mathrm{Id}+W)}\,dv\,V.
\]
The exponential inside the integral satisfies
$\lVert e^{-v(\mathrm{Id}+W)}\rVert_\infty
\le e^{-v}e^{v\lVert W\rVert_\infty}=1$.
Hence, for every $s\ge0$,
\begin{equation}\label{eq:average-semigroup}
 \lVert E_0(s)\rVert_\infty\le1+2s,
 \qquad \lVert E_0(s)\rVert_2\le1.
\end{equation}
The second inequality follows from the positive semidefiniteness of
$P$.  Likewise, $A/D=P+R$ is symmetric with smallest eigenvalue at
least $-\lVert R\rVert_2$, so
\begin{equation}\label{eq:full-semigroup-l2}
 E(s):=e^{-sA/D},\qquad
 \lVert E(s)\rVert_2\le e^{s\lVert R\rVert_2}.
\end{equation}

\paragraph{A finite perturbation expansion.}
It remains to transfer the maximum-norm estimate from $E_0$ to $E$.
We use a finite form of the bounded-perturbation expansion for semigroups
\cite[Theorem~III.1.10]{EngelNagel2000}, estimating its initial terms
and remainder in different norms.
Choose a fixed integer $K$ with $K\gamma>2$.  Iterating the
variation-of-constants identity
\[
 E(s)=E_0(s)-\int_0^s E_0(s-v)R E(v)\,dv
\]
gives terms with $j=0,\ldots,K-1$ copies of $R$, followed by a
remainder with $K$ copies.  The term with $j$ copies is integrated
over an ordered simplex of volume $s^j/j!$ and contains $j+1$
factors $E_0$, all at nonnegative times at most $s$.  By
\eqref{eq:average-semigroup}, its maximum norm is at most
\[
 (1+2s)^{j+1}\lVert R\rVert_\infty^j\frac{s^j}{j!}.
\]
The remainder has the exact form
\[
 \begin{split}
 \mathcal R_K(s)=(-1)^K
 \int_{0\le t_K\le\cdots\le t_1\le s}
 &E_0(s-t_1)R E_0(t_1-t_2)R\cdots\\[-2pt]
 &\hspace{18mm}\cdots E_0(t_{K-1}-t_K)R E(t_K)\,
 dt_K\cdots dt_1.
 \end{split}
\]
Its Euclidean norm is at most
\[
 \lVert\mathcal R_K(s)\rVert_2
 \le\lVert R\rVert_2^K e^{s\lVert R\rVert_2}\frac{s^K}{K!}.
\]
For an $m\times m$ matrix $H$, Cauchy--Schwarz gives
$\lVert H\rVert_\infty\le\sqrt m\,\lVert H\rVert_2$.
Here $m=|\cE|\le n^2/2$, so, uniformly for $s\le\log n$,
\[
 \lVert\mathcal R_K(s)\rVert_\infty
 \le n(Cn^{-\gamma})^K
       e^{Cn^{-\gamma}\log n}\frac{(\log n)^K}{K!}
 =o(1).
\]
The other $K$ terms have maximum norm bounded by a fixed power of
$1+\log n$, by~\eqref{eq:R-norms}.  This proves
\eqref{eq:semigroup}.  The order $K$ may be large, but is fixed
independently of $n$ and $G$.
\end{proof}

\subsection{Closing the probability estimate}

The semigroup bound now controls the cumulative error in
\eqref{eq:linearized-error}.  This is the only use of the spectral
calculation in the probability argument.

\begin{proposition}\label{prop:unfolding}
For every graph $G$ satisfying Proposition~\ref{prop:prefix}, the
independent-unfolding probabilities obey
\begin{equation}\label{eq:unfolding-asymptotics}
 q_e(t)=(1+o(1))q(t),\qquad
 q_{e,T}(t)=(1+o(1))q(t),\qquad
 I_{e,T}(t)=O\!\left(\frac{\log n}{D}\right).
\end{equation}
The estimates hold uniformly over $G$, $0\le t\le1$, and every
incident pair $e\in T\in\cT$; the first estimate also holds for every
$e\in\cE$.
\end{proposition}

\begin{proof}
Use $a$ and $\delta_n$ from~\eqref{eq:bootstrap-exponent}.
All $I_{e,T}$ vanish and all probabilities equal one at zero, so by
continuity the conditions~\eqref{eq:bootstrap-assumptions} hold on
some interval starting at zero.  On every such interval,
\eqref{eq:I-bootstrap} improves the second condition to
$I_{e,T}\le1/4$ for large $n$.

Write $h(s)=z'(s)+(A/D)z(s)$.  As $z(0)=0$, variation of constants
and Proposition~\ref{prop:semigroup} give, on the same interval,
\begin{align*}
 \lVert z(s)\rVert_\infty
 &\le\int_0^s
       \lVert e^{-(s-v)A/D}\rVert_\infty
       \lVert h(v)\rVert_\infty\,dv\\
 &\le C(1+\log n)^{J+1}
       \left(n^{-2a}+n^{-c}+\frac{\log n}{D}\right)
 =o(n^{-a}).
\end{align*}
The last equality follows from $a>0$, $c-a>0$ and
$2\epsilon-a>0$, together with $D\sim n^{2\epsilon}$: every
fixed power of $\log n$ is dominated by each of these positive powers
of $n$.  In particular, for large $n$ this improves the first
condition to $\lVert z\rVert_\infty\le\delta_n/2$.  The spectral
exponent $\gamma$ only fixes the integer $K$ in the semigroup estimate.

If either bootstrap condition had a first exit before $s=S$, these
strict improvements would hold up to that exit.  All product factors
would then remain at least $3/4$, so the finitely many functions and
their logarithms would be continuous on a neighborhood of the exit.
Both conditions would persist beyond it, a contradiction.  Thus the
bounds hold for $s<S$, and continuity gives them also at $s=S$;
the endpoint factors remain positive.  This proves
$z_e(s)=o(n^{-a})$ throughout $[0,S]$ and the asserted bound on
$I_{e,T}$.  Finally,
\[
 \frac{q_e(t)}{q(t)}=e^{z_e(s)},\qquad
 \frac{q_{e,T}(t)}{q(t)}
 =\frac{e^{z_e(s)}}{1-I_{e,T}(t)},
\]
which establishes both probability estimates.
\end{proof}

One consequence records the precision needed for the later exploration
of selected paths.  Set
$J_n=\max_{e\in T\in\cT,\,0\le t\le1}I_{e,T}(t)$.
At a paired call with focus $\{f,g\}\subset T$ and omitted parent
$T$, the product of all no-success factors is
$q_{f,T}(t)q_{g,T}(t)$.  Deleting at most two factors from this
product increases it by at most $(1-J_n)^{-2}$.  Hence, uniformly
in the call and threshold, for large $n$ the resulting product is at
most
\begin{equation}\label{eq:off-pattern-bound}
 \frac{q_{f,T}(t)q_{g,T}(t)}{(1-J_n)^2}
 \le 2q(t)^2=\frac2{1+2Dt}.
\end{equation}
This remains a statement about the independent unfolding.  We next
compare its queries with the finite priority process.
